\subsubsection{Satz des Pythagoras}
\label{subsubsec:fundamentals_right_triangles_pythagorean_theorem}

Der Satz des Pythagoras verknüpft die drei Seitenlängen eines rechtwinkligen
Dreiecks. Dabei bezeichnet $c$ die Hypotenuse und $a,b$ die Katheten.

\begin{TheoremBox}{Satz des Pythagoras}{fundamentals_pythagorean_theorem}
  \noindent
  \begin{minipage}[t]{0.64\linewidth}
    Für jedes Dreieck sind die folgenden Aussagen äquivalent:
    \begin{ClozeArea}[22]{3}
      \[
        \begin{gathered}
          \text{Dreieck ist rechtwinklig}\\[-0.1em]
          \gamma=90^\circ
        \end{gathered}
        \quad \Longleftrightarrow \quad
        a^2+b^2=c^2.
      \]
    \end{ClozeArea}
    Dabei ist die Seite $c$ die dem Winkel $\gamma$ gegenüberliegende Seite.
  \end{minipage}\hfill
  \begin{minipage}[t]{0.31\linewidth}
    \vspace{0pt}
    \centering
    \CachedTikzFigure{content/resources/chapters/fundamentals/tikz/right_triangle_pythagorean_example.tex}
  \end{minipage}
\end{TheoremBox}
Aufgrund der Äquivalenz beider Aussagen ist der Satz des Pythagoras nicht nur ein Werkzeug
um Seitenlängen in rechtwinkligen Dreiecken zu berechnen, sondern auch um die
Rechtwinkligkeit eines Dreiecks zu prüfen.

\begin{ExampleBox}{Satz des Pythagoras}{fundamentals_pythagorean_theorem_example}
  \begin{minipage}[c]{0.62\linewidth}
  Gegeben sind die Hypotenuse $c=5\,\mathrm{cm}$ und die Kathete
  $b=4\,\mathrm{cm}$. Gesucht ist die Kathete $a$.
    \begin{ClozeArea}[22]{7}
      \begin{align*}
        \phantom{\Longrightarrow \quad} c^2&=a^2+b^2,\\
        \Longrightarrow \quad a&=\sqrt{c^2-b^2}=\sqrt{(5\,\mathrm{cm})^2-(4\,\mathrm{cm})^2}
         =\sqrt{9\,\mathrm{cm}^2}=3\,\mathrm{cm}.
      \end{align*}
    \end{ClozeArea}
  \end{minipage}\hfill
  \begin{minipage}[c]{0.33\linewidth}
    \centering
    \CachedTikzFigure{content/resources/chapters/fundamentals/tikz/right_triangle_pythagorean_example.tex}
  \end{minipage}
\end{ExampleBox}
